% !TEX program = xelatex
\PassOptionsToPackage{fontset=none}{ctex}
\documentclass[aspectratio=169,10.5pt,xcolor={dvipsnames,svgnames}]{beamer}

\usetheme{Madrid}
\usecolortheme{default}
\setbeamertemplate{navigation symbols}{}

\usepackage[UTF8,scheme=plain]{ctex}
\usepackage{fontspec}
\usepackage{amsmath,amssymb,bm}
\usepackage{graphicx}
\usepackage{booktabs,array,tabularx}
\usepackage{multirow}
\usepackage{tikz}
\usetikzlibrary{positioning,arrows.meta,shapes.geometric,calc}
\usepackage{pgfplots}
\pgfplotsset{compat=1.18}
\usepackage{hyperref}
\usepackage{appendixnumberbeamer}

\IfFontExistsTF{PingFang SC}{\setCJKmainfont{PingFang SC}}{\IfFontExistsTF{Songti SC}{\setCJKmainfont{Songti SC}}{\setCJKmainfont{FandolSong-Regular}}}
\IfFontExistsTF{PingFang SC}{\setCJKsansfont{PingFang SC}}{\IfFontExistsTF{Heiti SC}{\setCJKsansfont{Heiti SC}}{\setCJKsansfont{FandolHei-Regular}}}
\IfFontExistsTF{Kaiti SC}{\setCJKfamilyfont{zhkai}{Kaiti SC}}{}

\definecolor{CsuBlue}{RGB}{120,20,48}
\definecolor{CsuTeal}{RGB}{143,53,78}
\definecolor{CsuGold}{RGB}{191,150,77}
\definecolor{SoftGray}{RGB}{88,88,88}
\definecolor{PanelGray}{RGB}{245,247,250}

\setbeamercolor{palette primary}{bg=CsuBlue,fg=white}
\setbeamercolor{palette secondary}{bg=CsuBlue!85,fg=white}
\setbeamercolor{palette tertiary}{bg=CsuBlue!75,fg=white}
\setbeamercolor{palette quaternary}{bg=CsuBlue!70,fg=white}
\setbeamercolor{structure}{fg=CsuBlue}
\setbeamercolor{frametitle}{bg=CsuBlue,fg=white}
\setbeamercolor{title}{fg=CsuBlue}
\setbeamercolor{subtitle}{fg=CsuTeal}
\setbeamercolor{normal text}{fg=black,bg=white}
\setbeamercolor{block title}{bg=CsuBlue!92,fg=white}
\setbeamercolor{block body}{bg=PanelGray,fg=black}
\setbeamercolor{block title example}{bg=CsuTeal!92,fg=white}
\setbeamercolor{block body example}{bg=CsuTeal!7,fg=black}
\setbeamercolor{block title alerted}{bg=CsuGold!92,fg=white}
\setbeamercolor{block body alerted}{bg=CsuGold!10,fg=black}
\setbeamercolor{item projected}{bg=CsuBlue,fg=white}

\setbeamerfont{title}{size=\LARGE,series=\bfseries}
\setbeamerfont{subtitle}{size=\large,series=\bfseries}
\setbeamerfont{frametitle}{size=\Large,series=\bfseries}
\setbeamerfont{author}{size=\normalsize}
\setbeamerfont{institute}{size=\small}
\setbeamerfont{block title}{size=\normalsize,series=\bfseries}
\setbeamerfont{normal text}{size=\small}

\setbeamertemplate{itemize item}{\color{CsuBlue}\large\textbullet}
\setbeamertemplate{itemize subitem}{\color{CsuTeal}\small\textbullet}
\setlength{\leftmargini}{1.3em}
\setlength{\leftmarginii}{1.1em}
\setlength{\parskip}{0.25em}
\linespread{1.08}

\newcommand{\hl}[1]{\textcolor{CsuGold!80!black}{\textbf{#1}}}
\newcommand{\hlb}[1]{\textcolor{CsuBlue}{\textbf{#1}}}
\newcommand{\hlt}[1]{\textcolor{CsuTeal}{\textbf{#1}}}
\newcommand{\metric}[2]{\textbf{#1}\,{\small #2}}

\title[绿色金融支持碳中和]{绿色金融支持碳中和的政策工具与效果评估}
\subtitle{北京师范大学本科毕业论文答辩}
\author[吴坤泽]{答辩人：吴坤泽，指导教师：林卫斌 教授}
\institute[BNU]{北京师范大学经济与工商管理学院，金融学}
\date{2026年5月}

\begin{document}

\begin{frame}[plain,noframenumbering]
  \begin{tikzpicture}[remember picture,overlay]
    \fill[CsuBlue] (current page.north west) rectangle ([yshift=-2.15cm]current page.north east);
    \fill[CsuGold] ([yshift=-2.15cm]current page.north west) rectangle ([yshift=-2.3cm]current page.north east);
  \end{tikzpicture}
  \vspace{0.8cm}

  {\usebeamerfont{title}\color{CsuBlue}\inserttitle\par}
  \vspace{0.4cm}
  {\usebeamerfont{subtitle}\color{CsuTeal}\insertsubtitle\par}
  \vspace{1.0cm}

  \begin{block}{研究聚焦}
    基于 2013--2022 年中国 31 个省级行政区面板数据，
    用 \hlb{双向固定效应模型} 评估绿色金融是否显著降低 \hl{碳排放强度}，
    并检验其区域异质性与政策含义。
  \end{block}

  \vfill
  \begin{tabular}{@{}ll@{}}
    答辩人： & 吴坤泽 \\
    指导教师： & 林卫斌 \\
    学院专业： & 经济与工商管理学院，金融学 \\
    学校： & 北京师范大学
  \end{tabular}
  \vspace{0.5cm}
\end{frame}

\begin{frame}{汇报提纲}
  \tableofcontents[hideallsubsections]
\end{frame}

\section{研究背景}

\begin{frame}{研究背景：双碳目标需要金融体系提供持续支持}
  \begin{columns}[T,onlytextwidth]
    \column{0.56\textwidth}
    \begin{itemize}
      \item 气候变化治理已从环境议题上升为 \hlb{经济结构转型议题}。
      \item 中国提出 \hl{2030 碳达峰、2060 碳中和}，低碳转型面临巨大资金缺口。
      \item 绿色金融通过信贷、债券、基金、保险等工具，引导资本流向绿色产业。
      \item 关键问题不只是“规模有没有增长”，而是 \hlt{是否真的带来了减排效果}。
    \end{itemize}

    \begin{block}{本文关注}
      绿色金融发展是否能够显著降低各省 \hlb{碳排放强度}，
      以及这种效果在东中西部地区是否存在明显差异。
    \end{block}

    \column{0.4\textwidth}
    \begin{exampleblock}{现实动机}
      \begin{itemize}
        \item 地区产业结构差异大
        \item 金融市场发育程度不同
        \item 政策资源投放效果可能不均衡
      \end{itemize}
    \end{exampleblock}
  \end{columns}
\end{frame}

\begin{frame}{研究意义：回答“绿色金融是否有效、何处更有效”}
  \begin{enumerate}
    \item 从学术上，为 \hlb{绿色金融与环境经济学交叉研究} 提供省级面板实证证据。
    \item 从政策上，检验绿色金融发展是否真正转化为 \hl{碳减排绩效}。
    \item 从实践上，识别东中西部效果差异，为 \hlt{差异化政策设计} 提供依据。
  \end{enumerate}

  \vspace{0.6em}
  \begin{alertblock}{答辩主线}
    政策背景提出问题，理论机制建立假设，实证结果检验效果，最后回到政策工具优化。
  \end{alertblock}
\end{frame}

\section{理论与假设}

\begin{frame}{理论基础：三条逻辑共同支撑研究框架}
  \begin{columns}[T,onlytextwidth]
    \column{0.32\textwidth}
    \begin{block}{外部性理论}
      碳排放具有典型负外部性，市场会对高碳活动 \hl{定价不足}。
    \end{block}

    \column{0.32\textwidth}
    \begin{exampleblock}{金融功能理论}
      金融体系通过 \hlb{资源配置、风险管理、信息筛选} 改变资金流向。
    \end{exampleblock}

    \column{0.32\textwidth}
    \begin{alertblock}{波特假说}
      合理环境约束与金融激励可推动企业技术升级与绿色创新。
    \end{alertblock}
  \end{columns}

  \vspace{0.8em}
  \begin{itemize}
    \item 理论上，绿色金融应通过融资约束与资源倾斜，降低高碳活动比重。
    \item 但不同地区资源禀赋不同，实际效果可能存在显著 \hl{区域异质性}。
  \end{itemize}
\end{frame}

\begin{frame}{研究假设与分析思路}
  \begin{block}{假设 H1}
    绿色金融发展水平提升，将显著降低各地区碳排放强度。
  \end{block}

  \begin{exampleblock}{假设 H2}
    绿色金融的减排效应具有区域异质性，在不同地区的边际作用不同。
  \end{exampleblock}

  \vspace{0.6em}
  \begin{itemize}
    \item 如果 H1 成立，说明绿色金融不是概念性政策，而是有可测量的减排效果。
    \item 如果 H2 成立，说明政策设计不能“一刀切”，而应强调 \hlb{区域适配性}。
  \end{itemize}
\end{frame}

\section{研究设计}

\begin{frame}{数据与模型设计}
  \begin{columns}[T,onlytextwidth]
    \column{0.50\textwidth}
    \begin{block}{样本与变量}
      \begin{itemize}
        \item 样本：2013--2022 年中国 31 个省级行政区
        \item 被解释变量：\hlb{碳排放强度}（对数）
        \item 核心解释变量：\hlb{绿色金融发展指数}（对数）
        \item 控制变量：人均 GDP、第二产业占比、煤炭消费占比等
      \end{itemize}
    \end{block}

    \column{0.46\textwidth}
    \begin{exampleblock}{模型方法}
      \begin{itemize}
        \item 双向固定效应模型
        \item 控制省份固定效应
        \item 控制年份固定效应
        \item 后续加入工具变量与稳健性检验
      \end{itemize}
    \end{exampleblock}
  \end{columns}
\end{frame}

\begin{frame}{识别策略：从基准回归到稳健性与异质性}
  \centering
  \begin{tikzpicture}[
    node distance=0.55cm,
    box/.style={draw=CsuBlue,rounded corners,fill=CsuBlue!6,minimum width=2.55cm,minimum height=0.95cm,align=center,font=\small},
    arr/.style={-{Stealth[length=2.2mm]},thick,CsuBlue}
  ]
    \node[box] (base) {基准回归\\双向固定效应};
    \node[box,right=of base] (iv) {内生性处理\\工具变量法};
    \node[box,right=of iv] (robust) {稳健性检验\\三种替换方案};
    \node[box,right=of robust] (hetero) {异质性分析\\东中西部分组};

    \draw[arr] (base) -- (iv);
    \draw[arr] (iv) -- (robust);
    \draw[arr] (robust) -- (hetero);
  \end{tikzpicture}

  \vspace{0.8em}
  \begin{itemize}
    \item 工具变量：\hl{滞后一期绿色金融发展指数}
    \item 稳健性：滞后一期、替换碳排放总量、剔除直辖市
    \item 异质性：比较东部、中部、西部地区的减排效果差异
  \end{itemize}
\end{frame}

\section{实证结果}

\begin{frame}{基准回归结果：绿色金融对碳排放强度有显著抑制作用}
  \begin{table}
    \centering
    \small
    \begin{tabular}{@{}lccc@{}}
      \toprule
      模型列 & (1) & (2) & (3) \\
      \midrule
      绿色金融发展指数 & 负 & 负 & \textbf{-2.363*} \\
      经济发展水平控制 & 否 & 是 & 是 \\
      全部控制变量 & 否 & 否 & 是 \\
      固定效应 & 双向 & 双向 & 双向 \\
      \bottomrule
    \end{tabular}
  \end{table}

  \begin{alertblock}{核心发现}
    在纳入全部控制变量后，绿色金融发展指数系数为 \hl{-2.363}，`p=0.057`，
    表明绿色金融发展水平每提高 1\%，碳排放强度平均下降约 \hl{2.36\%}。
  \end{alertblock}

  \begin{itemize}
    \item 人均 GDP 系数显著为负，说明经济发展水平提升有助于降低碳排放强度。
    \item 第二产业占比与煤炭消费占比方向符合预期，但在省级层面统计上不够稳定。
  \end{itemize}
\end{frame}

\begin{frame}{内生性与稳健性：核心结论保持稳定}
  \begin{columns}[T,onlytextwidth]
    \column{0.48\textwidth}
    \begin{block}{工具变量检验}
      \begin{itemize}
        \item 工具变量：滞后一期绿色金融发展指数
        \item 第一阶段 F 统计量约为 \hlb{$2.2\times10^4$}
        \item 第二阶段系数：\hl{-2.307***}
      \end{itemize}
    \end{block}

    \column{0.48\textwidth}
    \begin{exampleblock}{稳健性检验}
      \begin{itemize}
        \item 滞后一期解释变量：\hl{-2.242*}
        \item 替换为碳排放总量：\hl{-1.035**}
        \item 剔除直辖市样本：\hl{-3.668***}
      \end{itemize}
    \end{exampleblock}
  \end{columns}

  \begin{alertblock}{结论}
    无论从内生性处理还是多种稳健性检验看，绿色金融对碳排放的抑制效应都没有消失。
  \end{alertblock}
\end{frame}

\section{区域异质性}

\begin{frame}{异质性分析：减排效果主要体现在西部地区}
  \begin{table}
    \centering
    \small
    \begin{tabular}{@{}lccc@{}}
      \toprule
      地区 & 东部 & 中部 & 西部 \\
      \midrule
      绿色金融发展指数系数 & 不显著 & 不显著 & \textbf{-6.568*} \\
      减排效果判断 & 弱 & 弱 & 强 \\
      \bottomrule
    \end{tabular}
  \end{table}

  \begin{columns}[T,onlytextwidth]
    \column{0.48\textwidth}
    \begin{block}{西部地区为何更显著}
      \begin{itemize}
        \item 碳排放基数更高
        \item 边际减排空间更大
        \item 能源结构偏煤，改进空间明显
      \end{itemize}
    \end{block}

    \column{0.48\textwidth}
    \begin{exampleblock}{东中部为何不显著}
      \begin{itemize}
        \item 碳强度已相对较低
        \item 边际减排成本更高
        \item 绿色金融资源更多分散到非工业领域
      \end{itemize}
    \end{exampleblock}
  \end{columns}
\end{frame}

\begin{frame}{如何理解区域差异}
  \begin{itemize}
    \item 西部地区同时具备 \hlb{高排放基数 + 清洁能源项目集中布局} 两个条件。
    \item “西电东送”、大型风电光伏基地建设，使绿色金融资源更容易形成直接减排效果。
    \item 东中部地区产业结构更成熟，未来要释放绿色金融效能，关键在于 \hl{与产业升级、技术创新深度协同}。
  \end{itemize}

  \begin{alertblock}{研究含义}
    绿色金融政策有效，但其效果高度依赖地区的产业结构、能源禀赋与政策承接能力。
  \end{alertblock}
\end{frame}

\section{政策启示}

\begin{frame}{政策建议：从“做大总量”走向“精准施策”}
  \begin{enumerate}
    \item \hlb{持续完善绿色金融体系}：扩大绿色信贷、绿色债券、绿色保险、绿色基金等供给规模。
    \item \hlb{实施差异化区域政策}：西部继续聚焦清洁能源与节能技改，东中部强化绿色技术创新与产业升级。
    \item \hlb{强化政策协同}：推动绿色金融与产业政策、环境规制、碳交易体系联动。
    \item \hlb{完善基础设施}：统一绿色标准、建设信息共享平台、加强绿色金融人才培养。
  \end{enumerate}
\end{frame}

\begin{frame}{实践价值：对政府、金融机构与企业分别意味着什么}
  \begin{columns}[T,onlytextwidth]
    \column{0.32\textwidth}
    \begin{block}{政府}
      更需要看 \hl{政策效果}，而非只看工具投放规模。
    \end{block}

    \column{0.32\textwidth}
    \begin{exampleblock}{金融机构}
      需要降低“漂绿”风险，把资源投向 \hlb{真正有减排效果} 的项目。
    \end{exampleblock}

    \column{0.32\textwidth}
    \begin{alertblock}{企业}
      绿色融资不只是拿资金，更是推动低碳转型和竞争力重塑的契机。
    \end{alertblock}
  \end{columns}
\end{frame}

\section{结论展望}

\begin{frame}{研究结论}
  \begin{enumerate}
    \item 绿色金融发展对碳排放强度具有显著抑制作用，基准结果为 \hl{-2.363*}。
    \item 工具变量与多种稳健性检验均支持该结论，说明结论具有较强可靠性。
    \item 绿色金融减排效应存在明显区域异质性，\hlb{西部显著、东中部不显著}。
    \item 因此，绿色金融政策设计应从统一扩张转向 \hlt{区域差异化配置}。
  \end{enumerate}
\end{frame}

\begin{frame}{局限与未来研究方向}
  \begin{columns}[T,onlytextwidth]
    \column{0.48\textwidth}
    \begin{block}{当前局限}
      \begin{itemize}
        \item 绿色金融使用综合指数，难区分不同工具效应
        \item 样本停留在省级层面，时间跨度有限
        \item 因果识别仍有进一步加强空间
      \end{itemize}
    \end{block}

    \column{0.48\textwidth}
    \begin{exampleblock}{未来方向}
      \begin{itemize}
        \item 下沉到地级市或企业层面
        \item 细分绿色信贷、债券、保险等工具
        \item 引入双重差分等更强识别策略
      \end{itemize}
    \end{exampleblock}
  \end{columns}

  \vspace{0.6em}
  \begin{alertblock}{一句话总结}
    本文证明了绿色金融不仅是政策倡议，更是能够被实证识别的碳减排工具，但要真正发挥作用，关键在于 \hl{分地区、分工具、分场景优化配置}。
  \end{alertblock}
\end{frame}

\section{答辩致谢}

\begin{frame}[plain]
  \centering
  \vspace{1.8cm}
  {\usebeamerfont{title}\color{CsuBlue}谢谢各位老师聆听\par}
  \vspace{0.6cm}
  {\usebeamerfont{subtitle}\color{CsuTeal}敬请批评指正\par}
  \vspace{1.2cm}
  \begin{tabular}{@{}ll@{}}
    论文题目： & 绿色金融支持碳中和的政策工具与效果评估 \\
    答辩人： & 吴坤泽 \\
    学校： & 北京师范大学 \\
    学院： & 经济与工商管理学院
  \end{tabular}
\end{frame}


\end{document}
